\documentclass[12pt]{article}

\usepackage{amsmath}
\usepackage[margin=1in]{geometry}
\usepackage{fancyhdr}
\usepackage{amsthm}
\usepackage{amssymb}
\usepackage{pgfplots}
\usepackage{tikz}
\usepackage{bbm}
\usepackage{bm}
\usepackage{dsfont}
\usetikzlibrary{shapes.geometric}
\usepackage{enumerate}
\usepackage[shortlabels]{enumitem}

\pgfplotsset{compat=1.18}

\newcommand{\pderiv}[2][]{\frac{\partial#1}{\partial#2}}

\pagestyle{fancy}
\fancyhead[l]{Hudson Benites}
\fancyhead[c]{Day 9}
\fancyfoot[c]{\thepage}
\renewcommand{\headrulewidth}{0.2pt}
\setlength{\headheight}{15pt}
\fancyhead[r]{\today}

\newcommand{\magnitude}[1]{\left\lvert #1 \right\rvert} %chktex 1

\begin{document}
\begin{enumerate}[label=\Alph*., leftmargin=*]
    \item \begin{proof}
        Let $\beta$ be regular. We have that $\alpha=x\circ \beta$, and so $\dot{\alpha}=\nabla x \cdot \dot{\beta}=x_u\dot{\beta_1}+x_v\dot{\beta_2}$.
        We know that $\dot{\beta_1}$ and $\dot{\beta_2}$ never equal 0 at the same time. We also know that neither $x_u$ or $x_v$ are 0 because $x$ is a local parameterization.
        Therefore, $\dot{\alpha}$ is regular.

        Next, let $\alpha$ be regular. Similarly, $\dot{\alpha}=\nabla x \cdot \dot{\beta_1}=x_u\dot{\beta_2}+x_v\dot{\beta}$. $x_u$ and $x_v$ are linearly independent, and the LHS can never be 0, therefore $\dot{\beta_1}$ and $\dot{\beta_2}$ are
        never 0 at the same time. Therefore, $\beta$ is regular
    \end{proof}
    \item Below is my code to find the matrix $g$.\begin{figure}[h]\label{code}  
        \centering
            \includegraphics[width=10cm,keepaspectratio]{math.png}
            \caption{My code.}
        \end{figure}\par
        So, $g=\frac{4}{{(1+u^2+v^2)}^2}\mathds{1}$. This is a scalar of the identity matrix.
    \item \begin{enumerate}
        \item We compute that $x_u=(\cos v\cos u, \sin v\cos u, -\sin u)$. We also compute that $x_v=(-\sin u\sin v, \sin u\cos v, 0)$. Then,
        $x_u\cdot x_u=\cos^2 u\cos^2 v+\cos^2 u\sin^2 v + \sin^2 u=1$. Also, $x_u\cdot x_v=\cos v\cos u \sin u\sin v - \cos v\cos u \sin u\sin v = 0$. Finally,
        $x_v \cdot x_v = \sin^2 u\sin^2 v+\sin^2 u\cos^2 v = \sin^2 u$. So,\begin{equation*}
            g=\begin{bmatrix}
                1&0\\
                0&\sin^2 u
            \end{bmatrix}
        \end{equation*}.
        \item I could have used the fact that $J^{T}J=g$, where $J$ is the Jacobian.
        \item Let $x_1$ be the quasi-inverse of the stereographic project of $\mathbb{S}^2$ and $x_2$ be the parameterization of $\mathbb{S}^2$ using spherical coordinates.
        $x_1$ sends the point $(1,0)$ to $(1,0,0)$, and at that point $g_1=\frac{4}{{(1+1)}^2}\mathds{1}=\mathds{1}$. $x_2$ sends the point $(\frac{\pi}{2}, 0)$ to $(1,0,0)$, and $g_2=\begin{bmatrix}
            1&0\\
            0&\sin^2(\frac{\pi}{2})
        \end{bmatrix}=\mathds{1}$. So, they do have the same $g$ matrix at $(1,0,0)$.
        \item $x_1$ sends the point $(\sqrt{2}+1,0)$ to $(\frac{1}{\sqrt{2}},0,\frac{1}{\sqrt{2}})$. At that point $g_1=\frac{4}{{(2+\sqrt{2})}^2}\mathds{1}$. $x_2$ sends the point $(\frac{\pi}{4},0)$ to
        $(\frac{1}{\sqrt{2}},0,\frac{1}{\sqrt{2}})$. At that point $g_2=\begin{bmatrix}
            1&0\\
            0&\frac{1}{2}
        \end{bmatrix}$. By observation, $g_1\neq g_2$. So, $x_1$ and $x_2$ have different $g$ matrices at $(\frac{1}{\sqrt{2}},0,\frac{1}{\sqrt{2}})$.
    \end{enumerate}
\end{enumerate}
\end{document}